\subsection{Hörmander 乘子定理}
\label{sec:ch8-hormander-multiplier}

现在给出 Littlewood--Paley 理论的一些应用. 接下来三节讨论 $L^p$ 乘子的刻画问题. 更确切地说, 给定函数 $m$, 何时由
\[
 \widehat{T_mf}=m\widehat f
\]
定义的算子 $T_m$ 在 $L^p$, $1\le p\le\infty$, 上有界? 第~\ref{sec:ch3-multipliers}~节曾首先讨论 $L^p(\mathbb R)$ 上的乘子, 另见第~\ref{subsec:ch3-translation-invariant-operators}~节.

Sobolev 空间 $L_a^2(\mathbb R^n)$ 定义为所有满足 $(1+|\xi|^2)^{a/2}\widehat g(\xi)\in L^2$ 的函数 $g$, 后一函数的范数定义为 $g$ 在 $L_a^2$ 中的范数. 当 $a$ 是正整数时, 该定义等价于第~\ref{subsec:ch4-fractional-integrals-sobolev}~节给出的定义. 注意若 $a'<a$, 则 $L_a^2(\mathbb R^n)\subset L_{a'}^2(\mathbb R^n)$.

\begin{Proposition}
\label{prop:ch8-sobolev-l1}
若 $a>n/2$ 且 $g\in L_a^2(\mathbb R^n)$, 则 $\widehat g\in L^1$; 特别地, $g$ 连续且有界.
\end{Proposition}

\begin{Proof}
由 Cauchy--Schwarz 不等式,
\[
 \int_{\mathbb R^n}|\widehat g(\xi)|\,d\xi
 \le\left(\int_{\mathbb R^n}(1+|\xi|^2)^a|\widehat g(\xi)|^2\,d\xi\right)^{1/2}
 \left(\int_{\mathbb R^n}\frac{d\xi}{(1+|\xi|^2)^a}\right)^{1/2}
 \le C_a\|g\|_{L_a^2}.
\]
\end{Proof}

\hypertarget{chapter-08-page-176}{}
由此可知, 若 $m\in L_a^2$ 且 $a>n/2$, 则 $m$ 是 $L^p$, $1\le p\le\infty$, 上的乘子. 事实上, 若 $\widehat{Tf}=m\widehat f$, 则 $Tf=K*f$, 其中 $K\in L^1$. Hörmander 定理在弱得多的假设下仍保证 $m$ 是 $L^p$ 乘子. 证明前先建立一个加权范数不等式.

\begin{Lemma}
\label{lem:ch8-scaled-multiplier-weighted}
设 $m\in L_a^2$, $a>n/2$, $\lambda>0$. 定义
\[
 \widehat{T_\lambda f}(\xi)=m(\lambda\xi)\widehat f(\xi).
\]
则
\[
 \int_{\mathbb R^n}|T_\lambda f(x)|^2u(x)\,dx
 \le C\int_{\mathbb R^n}|f(x)|^2Mu(x)\,dx,
\]
其中常数 $C$ 与 $u,\lambda$ 无关, $M$ 是 Hardy--Littlewood 极大算子.
\end{Lemma}

\begin{Proof}
若 $\widehat K=m$, 则由假设 $R(x)=(1+|x|^2)^{a/2}K(x)\in L^2$, 而 $T_\lambda$ 的核为 $\lambda^{-n}K(\lambda^{-1}x)$. 因而
\[
\begin{aligned}
 \int_{\mathbb R^n}|T_\lambda f|^2u
 &=\int_{\mathbb R^n}\left|\int_{\mathbb R^n}
 \frac{\lambda^{-n}R(\lambda^{-1}(x-y))}
 {(1+|\lambda^{-1}(x-y)|^2)^{a/2}}f(y)\,dy\right|^2u(x)\,dx\\
 &\le\|m\|_{L_a^2}^2
 \int_{\mathbb R^n}\int_{\mathbb R^n}
 \frac{\lambda^{-n}|f(y)|^2}{(1+|\lambda^{-1}(x-y)|^2)^a}u(x)\,dy\,dx\\
 &\le C_a\|m\|_{L_a^2}^2\int_{\mathbb R^n}|f(y)|^2Mu(y)\,dy.
\end{aligned}
\]
第二个不等式使用 Cauchy--Schwarz 不等式和 $\|R\|_2=\|m\|_{L_a^2}$; 第三个不等式对 $x$ 积分并使用 $(1+|x|^2)^{-a}$ 在 $a>n/2$ 时径向递减且可积这一事实, 比较命题~\ref{prop:ch2-radial-majorization} 和定理~\ref{thm:ch2-weighted-maximal}.
\end{Proof}

为陈述 Hörmander 定理, 取径向函数 $\psi\in C^\infty$, 支于 $1/2\le|\xi|\le2$, 并满足
\[
 \sum_{j=-\infty}^{\infty}|\psi(2^{-j}\xi)|^2=1,
 \qquad\xi\ne0.
\]

\begin{Theorem}[Hörmander]
\label{thm:ch8-hormander-multiplier}
设对某个 $a>n/2$,
\[
 \sup_j\|m(2^j\cdot)\psi\|_{L_a^2}<\infty.
\]
则与乘子 $m$ 相伴的算子 $T$ 在 $L^p(\mathbb R^n)$, $1<p<\infty$, 上有界.
\end{Theorem}

\begin{Proof}
应用定理~\ref{thm:ch8-smooth-square-rn}. 先定义
\[
 \widehat{S_jf}(\xi)=\psi(2^{-j}\xi)\widehat f(\xi).
\]
由 $\psi$ 的选取, \eqref{eq:ch8-smooth-square-rn-lower} 成立. 再取支于 $1/4\le|\xi|\le4$ 且在 $1/2\le|\xi|\le2$ 上等于 $1$ 的 $C^\infty$ 函数 $\widetilde\psi$. 定义
\[
 \widehat{\widetilde S_jf}(\xi)=\widetilde\psi(2^{-j}\xi)\widehat f(\xi).
\]
则 $S_j\widetilde S_j=S_j$, 且 $\{\widetilde S_j\}$ 满足 \eqref{eq:ch8-smooth-square-rn}; 注意 $\widetilde\psi$ 未必满足 \eqref{eq:ch8-smooth-partition}. 因此
\[
 \|Tf\|_p
 \le C\left\|\left(\sum_j|S_jTf|^2\right)^{1/2}\right\|_p
 =C\left\|\left(\sum_j|S_jT\widetilde S_jf|^2\right)^{1/2}\right\|_p.
\]

\hypertarget{chapter-08-page-177}{}
$S_jT$ 的乘子为 $\psi(2^{-j}\xi)m(\xi)$. 由假设和引理~\ref{lem:ch8-scaled-multiplier-weighted},
\[
 \int_{\mathbb R^n}|S_jTf|^2u
 \le C\int_{\mathbb R^n}|f|^2Mu,
\]
其中 $C$ 与 $j$ 无关. 如定理~\ref{thm:ch8-weighted-sequence-square} 的证明, 由该不等式可对 $p>2$ 证明
\[
 \left\|\left(\sum_j|S_jTg_j|^2\right)^{1/2}\right\|_p
 \le C\left\|\left(\sum_j|g_j|^2\right)^{1/2}\right\|_p.
\]
结合 $g_j=\widetilde S_jf$ 的 \eqref{eq:ch8-smooth-square-rn}, 得
\[
 \|Tf\|_p
 \le C\left\|\left(\sum_j|\widetilde S_jf|^2\right)^{1/2}\right\|_p
 \le C\|f\|_p,
 \qquad p>2.
\]
当 $p<2$ 时由对偶性得到结论, 比较第~\ref{subsec:ch3-translation-invariant-operators}~节.
\end{Proof}

Hörmander 定理通常表述如下.

\begin{Corollary}
\label{cor:ch8-hormander-classical}
令 $k=[n/2]+1$. 若 $m$ 在原点之外属于 $C^k$, 且对 $|\beta|\le k$,
\begin{equation}
\label{eq:ch8-hormander-annular-condition}
 \sup_{R>0}R^{|\beta|}
 \left(\frac1{R^n}\int_{R<|\xi|<2R}|D^\beta m(\xi)|^2\,d\xi\right)^{1/2}
 <\infty,
\end{equation}
则 $m$ 是 $L^p$, $1<p<\infty$, 上的乘子. 特别地, 若
\[
 |D^\beta m(\xi)|\le C_\beta|\xi|^{-|\beta|},
 \qquad|\beta|\le k,
\]
则 $m$ 是乘子.
\end{Corollary}

\begin{Proof}
在 \eqref{eq:ch8-hormander-annular-condition} 中作变量替换 $\xi\mapsto R\xi$, 并用
\[
 D^\beta m(R\cdot)(\xi)=R^{|\beta|}(D^\beta m)(R\xi),
\]
得
\[
 \sup_R\left(\int_{1<|\xi|<2}|D^\beta m(R\cdot)(\xi)|^2\,d\xi\right)^{1/2}
 \le C.
\]

\hypertarget{chapter-08-page-178}{}
由于
\[
 D^\beta(m(2^j\cdot)\psi)(\xi)
 =\sum_{|\gamma|\le|\beta|}C_{\gamma,\beta}
 D^\gamma m(2^j\cdot)(\xi)D^{\beta-\gamma}\psi(\xi),
\]
且 $|D^\beta\psi|\le C$, 可得 $\sup_j\|m(2^j\cdot)\psi\|_{L_k^2}<\infty$, 因而可应用定理~\ref{thm:ch8-hormander-multiplier}.
\end{Proof}

\begin{Example}
\label{ex:ch8-hormander-multiplier-examples}
\begin{enumerate}[label=(\arabic*)]
\item $m(\xi)=|\xi|^{it}$ 满足推论~\ref{cor:ch8-hormander-classical} 的后一假设, 因为 $|D^\beta m(\xi)|\le C_t|\xi|^{-|\beta|}$.
\item 若 $m$ 是零次齐次函数, 且在单位球面上属于 $C^k$, 其中 $k>[n/2]$, 则 $m$ 是 $L^p(\mathbb R^n)$, $1<p<\infty$, 上的乘子.
\end{enumerate}
\end{Example}
